% Responsable: Persona 4. ~300 palabras.
\section{Introducción y contexto del problema}

\subsection{Por qué los datos geoespaciales necesitan coordenadas planas}

La latitud $\lat$ y la longitud $\lon$ son ángulos, no distancias. Un grado de
latitud mide siempre unos 111.2~km, pero uno de longitud mide
$R\cos\lat\,\pi/180$: 111.2~km en el ecuador y 97.3~km a $29^\circ$\,N, en
Sonora. Los algoritmos del cliente (k-means, búsqueda de vecinos cercanos con
árboles k-d, interpolación espacial por IDW o kriging, mapas de calor por
densidad de núcleo y optimización de rutas) suponen distancias euclidianas.
Con millones de registros lo práctico es proyectar una sola vez a $(x, y)$ y
trabajar en el plano, en lugar de calcular la distancia sobre la esfera para
cada par. La pregunta es cuánto error mete ese paso.

\subsection{Por qué no existe un mapa perfecto}

En el siglo XVIII Euler demostró que ningún mapa plano conserva todas las
distancias de la esfera \cite{snyder1987}. La razón de fondo es el Teorema
Egregium de Gauss:
la curvatura gaussiana no cambia bajo isometrías, y la de la esfera es $1/R^2$
mientras que la del plano es 0 \cite{docarmo1976}. Toda proyección renuncia a
algo: las conformes conservan los ángulos, las equivalentes conservan las
áreas y las equidistantes conservan las distancias solo en ciertas
direcciones.

\subsection{Qué significa que una proyección sea conforme}

Una proyección es conforme si conserva los ángulos o, lo que es lo mismo, si
en cada punto estira todas las direcciones por el mismo factor $\sigma$. Eso
es lo que le importa al cliente: una distancia corta en el mapa es $\sigma$
veces la real sin importar su orientación, así que vecinos y grupos solo
dependen de cuánto varía $\sigma$ dentro del país. Hoy el cliente usa
Mercator; INEGI usa la Cónica Conforme de Lambert con paralelos estándar en
$17.5^\circ$ y $29.5^\circ$\,N, registrada como EPSG:6372 \cite{epsg} (a veces
se cita como EPSG:6361, un código que no aparece en el registro como sistema
de coordenadas). Esta wiki
compara las dos con la proyección conforme óptima para México y responde
cuál conviene, cuánto error mete cada una y cuánto se puede ganar.

\subsection{Variables del problema}

\begin{itemize}
  \item $\lat$, $\lon$: latitud y longitud (grados en el texto, radianes en las
    fórmulas); $R = 6\,371.0088$~km, el radio medio de la Tierra.
  \item $q(\lat)$: latitud isométrica, y $w = \lon + i\,q$, la coordenada
    compleja sobre la esfera (sección~\ref{sec:fundamentos}).
  \item $z = x + iy = f(w)$: coordenadas planas, en metros.
  \item $\sigma$: factor de escala, la longitud en el mapa dividida entre la
    real. En Mercator y Lambert depende solo de $\lat$; en la proyección
    óptima depende también de la forma de la región.
  \item $\region$: México continental (de $14.53^\circ$ a $32.72^\circ$\,N y de
    $86.7^\circ$ a $117.1^\circ$\,O), con frontera $\frontera$, y
    $\dist$, la distorsión de una proyección sobre $\region$ (Paso~3).
\end{itemize}
